\documentclass[10pt,a4paper]{amsart}
\usepackage[margin=19mm]{geometry}
\usepackage{amsmath,amssymb,mathtools}
\usepackage[scr=rsfs,cal=euler]{mathalfa}
\usepackage{xfrac}
\usepackage[unicode,hidelinks,bookmarks=false]{hyperref}
\newcommand{\ie}{i.e.\ }
\newcommand{\cf}{cf.\ }
\newcommand{\resp}{respectively\ }
\newcommand{\Subsection}{\subsection}
\providecommand{\nobreakdash}{\nobreak\hbox{-}}
\setlength{\emergencystretch}{2em}
\multlinegap0pt
\usepackage{dsfont}
\usepackage{amssymb}
\usepackage{caption}
\newtheorem{assumption}{Assumption}
\newtheorem{prop}{Proposition}
\newtheorem{theorem}{Theorem}
\newtheorem{lemma}{Lemma}
\def\P{{\mathcal P}}
\def\argmin{\mathop{\rm arg\, min}}
\title{Reinforcement Learning with Function Approximation for Non-Markov Processes}
\author{Ali Devran Kara}
\date{Source: arXiv:2601.00151v1 (CC BY 4.0)}
\begin{document}
\maketitle

\noindent\emph{Source and licence.} Reinforcement Learning with Function Approximation for Non-Markov Processes. Version arXiv:2601.00151v1 (CC BY 4.0), distributed by arXiv under the Creative Commons Attribution 4.0 International licence, http://creativecommons.org/licenses/by/4.0/, as recorded in the arXiv licence metadata for this exact version and shown on the abstract page https://arxiv.org/abs/2601.00151v1. Full licence notice: \texttt{licenses/arxiv-2601.00151-CC-BY-4.0.txt}.\par\smallskip

\noindent\textit{Editorial scope.} Counted unit: the article's theorem main\_thm (source0.tex lines 651--653). The article's complete original proof of it is delivered with it (source0.tex lines 655--718, one \texttt{proof} environment of 64 lines). No mathematics is written, repaired or re-derived: the statement and the proof are the article's own lines, in the article's own notation, and no retained line is substituted or re-flowed.  Retained uncounted as the source-local context the proof needs: lines 313--330; lines 344--346; lines 419--434; lines 589--591; lines 615--617; lines 621--623; lines 627--629; lines 647--649.  Nothing was omitted from any retained span, so the package carries no omission record.  No retained line contains a citation, so the wrapper carries no bibliography and no bibliography entry.  No retained line contains a figure environment, an image-inclusion command or drawing code, so no image is delivered and the package declares no asset dependency.  Editorial formatting only, in addition: an AMS \texttt{amsart} preamble that restates the article's own declarations and macros used by the retained lines, namely the environments \texttt{assumption}, \texttt{prop}, \texttt{theorem}, \texttt{lemma} on separate counters, together with the standard packages those lines need.  The retained environments print in this document as Assumption 1, Proposition 1, Theorem 1, Lemma 1 in order of appearance, whereas the article prints the same items with the numbering of its own full document.  The complete native source of the article as distributed in the arXiv e-print for this exact version is bundled unchanged as \texttt{sources/191975/source0.tex}.
\begin{assumption}\label{poisson}
\begin{itemize}
\item[i.] For any bounded function $f$, we have
\begin{align*}
\frac{1}{N} \sum_{t=1}^N f(Z_t) \to \int f(z) \pi(dz) \quad \text{a.s.}
\end{align*}
almost surely for some probability measure $\pi \in \P(\mathds{S}^2 \times \mathds{R} \times \mathds{U})$.

\item[ii.] For the matrix-valued function $A(Z_t)$ and the vector-valued function $b(Z_t)$, define
\begin{align}\label{gordin}
Y^A_t := \sum_{k=0}^\infty \| E[A(Z_{t+k})|\mathcal{F}_t] - A \|, \quad
Y^b_t := \sum_{k=0}^\infty \| E[b(Z_{t+k})|\mathcal{F}_t] - b \|,
\end{align}
where $A := \int A(z) \pi(dz)$ and $b := \int b(z) \pi(dz)$ and where we use the spectral norm for the matrices. We assume that these sequences are uniformly bounded in $L_2$: $\sup_t \|Y^A_t\|_2 < \infty$ and $\sup_t \|Y^b_t\|_2 < \infty$.

\item[iii.] $A(Z_t)$ and $b(Z_t)$ are uniformly bounded functions.
\end{itemize}
\end{assumption}

\begin{align}\label{mixing}
\alpha(k):=\sup_j \alpha(\mathcal{F}_j^- , \mathcal{F}^+_{j+k}).
\end{align}

\begin{prop}\label{key_prop}
Suppose Assumption \ref{poisson} holds (or Assumption \ref{mixing} as a sufficient condition for Assumption \ref{poisson}) and that the stationary average matrix $A$ is positive definite. Consider the stochastic approximation iteration
\[
\theta_{t+1} = \theta_t + \alpha_t \big( -A(Z_t) \theta_t + b(Z_t) \big),
\]
where $A(Z_t)$ and $b(Z_t)$ are matrix and vector valued functions, respectively. Then, $\theta_t$ converges almost surely to a limit $\theta^*$ satisfying
$
A \theta^* = b,
$
where
\[
A = E[A(Z)] = \int A(z) \, \pi(dz), \quad 
b = E[b(Z)] = \int b(z) \, \pi(dz),
\]
and $\pi$ is the stationary distribution of the joint process $Z_t = \{S_{t+1}, S_t, C_t, U_t\}$.
\end{prop}

\begin{align}\label{appr_bell}
T^\gamma f(s) := \int_\mathds{U}\left(c(s,u) + \beta \int f(s_1)\eta(ds_1|s,u)\right)\gamma(du|s).
\end{align}

\begin{assumption}\label{lin_ind}
We assume  for the rest of the paper that $\{\phi^i(s)\}$ are linearly independent in $L_2(\pi)$ such that $E\left[\Phi(S)\Phi^\intercal(S)\right]$ is invertible. 
\end{assumption}

\begin{align}\label{proj}
\theta_f =\argmin_{\theta\in\mathds{R}^d} \sqrt{\int_{\mathds{S}} \left| f(s) - \theta^\intercal {\bf \Phi}(s)\right|^2 \pi(ds)}.
\end{align}

\begin{prop}\label{comp_contr}
The mapping $\Pi T^\gamma$ is a contraction under the $L_2$ norm, and thus admits a unique fixed point.
\end{prop}

\begin{align}\label{main_iter}
\theta_{t+1}= \theta_t - \alpha_t {\bf \Phi}(S_t) \left[\theta_t^\intercal {\bf \Phi}(S_t) - C_t - \beta \theta^\intercal_t{\bf \Phi}(S_{t+1}) \right]
\end{align}

\begin{theorem}\label{main_thm}
Under Assumption \ref{poisson} and \ref{lin_ind}, if the learning rates are such that $\sum_t\alpha_t = \infty$ and $\sum_t \alpha_t^2<\infty$, then the iterations in (\ref{main_iter}) converge to some $\theta^* \in\mathds{R}^d$. Denoting by $V(s):={\theta^*}^\intercal {\bf \Phi}(s)$, $V(s)$ is the fixed point of the joint mapping $\Pi T^\gamma $ where the mappings $\Pi$ and $T^\gamma$ are defined in (\ref{proj}) and (\ref{appr_bell}).
\end{theorem}

\begin{proof}
We use Proposition \ref{key_prop} with 
\begin{align*}
&A(S_t,S_{t+1})= -\beta{\bf \Phi}(S_t){\bf \Phi}^\intercal(S_{t+1})+ {\bf \Phi}(S_t){\bf \Phi}^\intercal(S_{t})\\
&b(S_t,C_t)={\bf \Phi}(S_t)C_t.
\end{align*}
The matrices $A$ and $b$ are defined under the invariant measure $\pi$ of the  joint process $(S_t,S_{t+1},U_t,C_t)$. 
 
 We need to show that the matrix $A$ is positive definite. 
 
 \begin{lemma}\label{helpful_lem}
\begin{align*}
(\theta-\theta^*)^\intercal E\left[ {\bf \Phi}(S)\left[ C+\beta \theta^\intercal{\bf \Phi}(S_1) - \theta^\intercal{\bf \Phi}(S)\right] \right]<0
\end{align*}
for any $\theta\neq\theta^*$ where $\theta^*$ corresponds to the fixed point of the operator $\Pi T^\gamma$, that is ${\theta^*}^\intercal{\bf \Phi}(s)$, where $\theta^*$ is unique under Assumption \ref{lin_ind}. 
\end{lemma}
\begin{proof}
Recall that $\Pi(C+\beta \theta^\intercal{\bf \Phi}(S_1) )$ denotes the projection map on the span of $\{\phi^i(s)\}$. Note that
\begin{align*}
\Pi(C+\beta \theta^\intercal{\bf \Phi}(S_1) ) &= \Pi\left(E\left[C+\beta \theta^\intercal{\bf \Phi}(S_1) |S\right]\right)\\
&=\Pi\left( T^\gamma ( \theta^\intercal{\bf \Phi}(S) ) \right).
\end{align*}
The first order conditions imply that $E\left[{\bf \Phi}(S)[ C+\beta \theta^\intercal{\bf \Phi}(S_1)  - \Pi(C+\beta \theta^\intercal{\bf \Phi}(S_1) ) ] \right]=0$. Then,  by adding and subtracting $ \Pi(C+\beta \theta^\intercal{\bf \Phi}(S_1) )$:
\begin{align*}
&(\theta-\theta^*)^\intercal E[ {\bf \Phi}(S)[ C+\beta \theta^\intercal{\bf \Phi}(S_1)  - \Pi(C+\beta \theta^\intercal{\bf \Phi}(S_1) )  \\
&\qquad\qquad\qquad +\left(  \Pi(C+\beta \theta^\intercal{\bf \Phi}(S_1) )-\theta^\intercal{\bf \Phi}(S)\right)] ]\\
& = (\theta-\theta^*)^\intercal E[ {\bf \Phi}(S)\left(  \Pi(C+\beta \theta^\intercal{\bf \Phi}(S_1) )-\theta^\intercal{\bf \Phi}(S) \right)] .
\end{align*}


In what follows, we use the equality $\Pi(C+\beta \theta^\intercal{\bf \Phi}(S_1) ) = \Pi\left( T^\gamma ( \theta^\intercal{\bf \Phi}(S) ) \right)$, and we add and subtract ${\theta^*}^\intercal{\bf \Phi}(S) =\Pi T^\gamma({\theta^*}^\intercal{\bf \Phi}(S)  )$ to use the contraction property of the composition operator $\Pi T^\gamma$ (see Proposition \ref{comp_contr}):
\begin{align*}
& (\theta-\theta^*)^\intercal E[ {\bf \Phi}(S)\left(  \Pi(C+\beta \theta^\intercal{\bf \Phi}(S_1) )-\theta^\intercal{\bf \Phi}(S) \right)] \\
&= (\theta-\theta^*)^\intercal E[ {\bf \Phi}(S) (\Pi T^\gamma(\theta^\intercal{\bf \Phi}(S))-  {\theta^*}^\intercal{\bf \Phi}(S)  ) ]\\
& +  (\theta-\theta^*)^\intercal E[ {\bf \Phi}(S) ( {\theta^*}^\intercal{\bf \Phi}(S) -  \theta^\intercal{\bf \Phi}(S)  )]\\
&\leq  \| (\theta-\theta^*)^\intercal  {\bf \Phi}(S)\|_2 \| \Pi T^\gamma(\theta^\intercal{\bf \Phi}(S))-  {\theta^*}^\intercal{\bf \Phi}(S)\|_2\\
& - \| (\theta-\theta^*)^\intercal  {\bf \Phi}(S)\|^2_2\\
& \leq  (\beta-1) \| (\theta-\theta^*)^\intercal  {\bf \Phi}(S)\|^2_2 <0
\end{align*}
where we used the Cauchy-Schwarz inequality, and the $L_2$ norm is with respect to the invariant measure $\pi$. The last step follows from the uniqueness of $\theta^*$.
\end{proof}
We then have that
 \begin{align*}
 &(-A) (\theta-\theta^*)=E\left[\beta {\bf \Phi}(S){\bf \Phi}^\intercal(S_{1})- {\bf \Phi}(S){\bf \Phi}^\intercal(S)  \right] (\theta-\theta^*)\\
 &= E\left[{\bf \Phi}(S)\left( C+\beta {\bf \Phi}^\intercal(S_{1})\theta- {\bf \Phi}^\intercal(S)\theta \right) \right]\\
 & \quad - E\left[{\bf \Phi}(S)\left( C+\beta {\bf \Phi}^\intercal(S_{1})\theta^*- {\bf \Phi}^\intercal(S)\theta^* \right) \right]\\
 &=E\left[{\bf \Phi}(S)\left( C+\beta {\bf \Phi}^\intercal(S_{1})\theta- {\bf \Phi}^\intercal(S)\theta \right) \right]
 \end{align*}
 where the last step follows from the fact that ${\bf \Phi}^\intercal(S)\theta^*$ is the fixed point of the operator $\Pi T^\gamma$ and that $\Pi(C+\beta \theta^\intercal{\bf \Phi}(S_1) ) = \Pi\left( T^\gamma ( \theta^\intercal{\bf \Phi}(S) ) \right)$.  Together with Lemma \ref{helpful_lem}, this shows that 
 \begin{align*}
 (\theta-\theta^*)(-A)(\theta-\theta^*)<0
 \end{align*}
 for all $\theta\neq\theta^*$, and thus using Proposition \ref{key_prop} we can conclude that $\theta_t$ converges to some $\theta'$ that satisfies $A\theta'=b$, which implies that
 \begin{align*}
 E\left[{\bf \Phi}(S)\left( C+\beta {\bf \Phi}^\intercal(S_{1})\theta'- {\bf \Phi}^\intercal(S)\theta' \right) \right]=0
 \end{align*}
 then as argued earlier, $\theta'$ also satisfies:
 \begin{align*}
 &E\bigg[{\bf \Phi}(S)\bigg( T^\gamma({\bf \Phi}^\intercal(S)\theta' )- {\bf \Phi}^\intercal(S)\theta' \bigg) \bigg]=0
 \end{align*}
 which in turn implies that ${\bf \Phi}^\intercal(S)\theta' $ is the fixed point of the operator $\Pi T^\gamma$. Since the fixed point is unique, we have that $\theta'=\theta^*$ which completes the proof.


\end{proof}

\end{document}
